\documentclass[10pt,aspectratio=169,serif]{beamer}

\usepackage[fontset=none, scheme=plain]{ctex}
\IfFontExistsTF{Songti SC}{%
	\setCJKmainfont{Songti SC}[AutoFakeSlant=0.2]%
	\setCJKsansfont{Songti SC}[AutoFakeSlant=0.2]%
}{%
	\setCJKmainfont{FandolSong-Regular}[Extension=.otf,BoldFont=FandolSong-Bold,AutoFakeSlant=0.2]%
	\setCJKsansfont{FandolSong-Regular}[Extension=.otf,BoldFont=FandolSong-Bold]%
}
\usepackage{amsmath,amssymb,graphicx,tikz,booktabs,multicol,calligra}
\usepackage{pgfplots}
\pgfplotsset{compat=1.18}
\usepackage{tcolorbox}
\tcbuselibrary{skins}
\usepackage{XyygU}
\usepackage{caption}
\captionsetup[figure]{name=图,labelsep=period}

\usetikzlibrary{arrows.meta,positioning,calc}

% 公式间距
\setlength{\abovedisplayskip}{0pt plus 2pt minus 2pt}
\setlength{\belowdisplayskip}{0pt plus 2pt minus 2pt}

\title[复习参考题1]{复习参考题1}
\subtitle{集合与常用逻辑用语}
\author{}
\institute{}
\date{2026-10-16}

\begin{document}

\begin{frame}
  \titlepage
\end{frame}

% ========== 第1题 ==========
\begin{frame}[t]{第1题：用列举法表示集合}
\textbf{1.} 用列举法表示下列集合：

\vspace{0.5em}
(1) $A = \{x \mid x^2 = 9\}$；

\vspace{0.3em}
(2) $B = \{x \in \mathbf{N} \mid 1 \le x \le 2\}$；

\vspace{0.3em}
(3) $C = \{x \mid x^2 - 3x + 2 = 0\}$。
\end{frame}

\begin{frame}[t]{第1题 解答}
\begin{tcolorbox}[colback=orange!10,colframe=orange,title=\textbf{解}]
(1) 由 $x^2 = 9$，得 $x = \pm 3$，所以
\[
A = \{-3,\ 3\}.
\]

\vspace{0.5em}
(2) 满足 $1 \le x \le 2$ 的自然数为 $1, 2$，所以
\[
B = \{1,\ 2\}.
\]

\vspace{0.5em}
(3) 由 $x^2 - 3x + 2 = 0$，因式分解得 $(x-1)(x-2) = 0$，所以 $x = 1$ 或 $x = 2$，
\[
C = \{1,\ 2\}.
\]
\end{tcolorbox}
\end{frame}

% ========== 第2题 ==========
\begin{frame}[t]{第2题：动点集合的图形}
\textbf{2.} 设 $P$ 表示平面内的动点，属于下列集合的点组成什么图形？

\vspace{0.5em}
(1) $\{P \mid PA = PB\}$（$A$，$B$ 是两个不同定点）；

\vspace{0.3em}
(2) $\{P \mid PO = 3\ \text{cm}\}$（$O$ 是定点）。
\end{frame}

\begin{frame}[t]{第2题 解答}
\begin{tcolorbox}[colback=orange!10,colframe=orange,title=\textbf{解}]
(1) 到两个定点 $A$、$B$ 距离相等的点的集合，是线段 $AB$ 的\textbf{垂直平分线}。

\vspace{0.8em}
(2) 到定点 $O$ 距离等于 $3$ cm 的点的集合，是以 $O$ 为圆心、$3$ cm 为半径的\textbf{圆}。
\end{tcolorbox}
\end{frame}

% ========== 第3题 ==========
\begin{frame}[t]{第3题：两个垂直平分线的交点}
\textbf{3.} 设平面内有 $\triangle ABC$，且 $P$ 表示这个平面内的动点，指出属于集合
\[
\{P \mid PA = PB\} \cap \{P \mid PA = PC\}
\]
的点是什么。
\end{frame}

\begin{frame}[t]{第3题 解答}
\begin{tcolorbox}[colback=orange!10,colframe=orange,title=\textbf{解}]
$\{P \mid PA = PB\}$ 是线段 $AB$ 的垂直平分线；

\vspace{0.3em}
$\{P \mid PA = PC\}$ 是线段 $AC$ 的垂直平分线。

\vspace{0.8em}
它们的交集是两条垂直平分线的\textbf{交点}，即 $\triangle ABC$ 的\textbf{外心}（三角形外接圆的圆心，即三边垂直平分线的交点）。

\vspace{0.3em}
这个点到 $A$、$B$、$C$ 三个顶点的距离相等。
\end{tcolorbox}
\end{frame}

% ========== 第4题 ==========
\begin{frame}[t]{第4题：充分必要条件填空}
\textbf{4.} 请用"充分不必要条件""必要不充分条件""充要条件""既不充分也不必要条件"填空：

\vspace{0.5em}
(1) 三角形两边上的高相等是这个三角形为等腰三角形的 \underline{\hspace{3cm}}；

\vspace{0.5em}
(2) $x \in A$ 是 $x \in A \cup B$ 的 \underline{\hspace{3cm}}；

\vspace{0.5em}
(3) $x \in A$ 是 $x \in A \cap B$ 的 \underline{\hspace{3cm}}；

\vspace{0.5em}
(4) $x$，$y$ 为无理数是 $x + y$ 为无理数的 \underline{\hspace{3cm}}。
\end{frame}

\begin{frame}[t]{第4题 解答}
\begin{tcolorbox}[colback=orange!10,colframe=orange,title=\textbf{解}]
(1) \textbf{充要条件。}

\vspace{0.3em}
\small（三角形面积 $S = \frac{1}{2} a h_a = \frac{1}{2} b h_b$，若 $h_a = h_b$，则 $a = b$，即等腰三角形；反之亦然。）

\vspace{0.6em}
(2) \textbf{充分不必要条件。}

\vspace{0.3em}
\small（$x \in A \Rightarrow x \in A \cup B$，但反之不成立，因为 $x$ 可能只属于 $B$。）

\vspace{0.6em}
(3) \textbf{必要不充分条件。}

\vspace{0.3em}
\small（$x \in A \cap B \Rightarrow x \in A$，但反之不成立，因为 $x$ 可能只属于 $A$ 而不属于 $B$。）

\vspace{0.6em}
(4) \textbf{既不充分也不必要条件。}

\vspace{0.3em}
\small（反例：$\sqrt{2} + (-\sqrt{2}) = 0$ 是有理数；而 $\sqrt{2} + \sqrt{3}$ 是无理数，但其中 $0$ 不是无理数。）
\end{tcolorbox}
\end{frame}

% ========== 第5题 ==========
\begin{frame}[t]{第5题：判断充分必要条件的真假}
\textbf{5.} 已知 $a$，$b$，$c$ 是实数，判断下列命题的真假：

\vspace{0.5em}
(1) "$a > b$"是"$a^2 > b^2$"的充分条件； \hfill ( \qquad )

\vspace{0.3em}
(2) "$a > b$"是"$a^2 > b^2$"的必要条件； \hfill ( \qquad )

\vspace{0.3em}
(3) "$a > b$"是"$ac^2 > bc^2$"的充分条件； \hfill ( \qquad )

\vspace{0.3em}
(4) "$a > b$"是"$ac^2 > bc^2$"的必要条件。 \hfill ( \qquad )
\end{frame}

\begin{frame}[t]{第5题 解答}
\begin{tcolorbox}[colback=orange!10,colframe=orange,title=\textbf{解}]
(1) \textbf{假。} 反例：$a = 1$，$b = -2$，则 $a > b$，但 $a^2 = 1 < 4 = b^2$。

\vspace{0.5em}
(2) \textbf{假。} 反例：$a = -2$，$b = 1$，则 $a^2 = 4 > 1 = b^2$，但 $a < b$。

\vspace{0.5em}
(3) \textbf{假。} 当 $c = 0$ 时，$ac^2 = bc^2 = 0$，即使 $a > b$，也没有 $ac^2 > bc^2$。

\vspace{0.5em}
(4) \textbf{真。} 由 $ac^2 > bc^2$ 可知 $c^2 > 0$（否则两边都为 $0$），所以两边同除以 $c^2$ 得 $a > b$。
\end{tcolorbox}
\end{frame}

% ========== 第6题 ==========
\begin{frame}[t]{第6题：用符号表示含量词的命题并判断真假}
\textbf{6.} 用符号"$\forall$"与"$\exists$"表示下列含有量词的命题，并判断真假：

\vspace{0.5em}
(1) 任意实数的平方大于或等于 $0$；

\vspace{0.3em}
(2) 对任意实数 $a$，二次函数 $y = x^2 + a$ 的图象关于 $y$ 轴对称；

\vspace{0.3em}
(3) 存在整数 $x$，$y$，使得 $2x + 4y = 3$；

\vspace{0.3em}
(4) 存在一个无理数，它的立方是有理数。
\end{frame}

\begin{frame}[t]{第6题 解答}
\begin{tcolorbox}[colback=orange!10,colframe=orange,title=\textbf{解}]
(1) $\forall x \in \mathbf{R}$，$x^2 \ge 0$。\textbf{真命题。}

\vspace{0.5em}
(2) $\forall a \in \mathbf{R}$，二次函数 $y = x^2 + a$ 的图象关于 $y$ 轴对称。\textbf{真命题。}

\vspace{0.3em}
\small（因为 $(-x)^2 + a = x^2 + a$，函数是偶函数。）

\vspace{0.5em}
(3) $\exists x, y \in \mathbf{Z}$，$2x + 4y = 3$。\textbf{假命题。}

\vspace{0.3em}
\small（左边 $2(x+2y)$ 是偶数，右边 $3$ 是奇数，不可能相等。）

\vspace{0.5em}
(4) $\exists x \in \{\text{无理数}\}$，$x^3 \in \mathbf{Q}$。\textbf{真命题。}

\vspace{0.3em}
\small（取 $x = \sqrt[3]{2}$，则 $x$ 是无理数，$x^3 = 2$ 是有理数。）
\end{tcolorbox}
\end{frame}

% ========== 第7题 ==========
\begin{frame}[t]{第7题：写命题的否定并判断真假}
\textbf{7.} 写出下列命题的否定，并判断它们的真假：

\vspace{0.5em}
(1) $\forall a \in \mathbf{R}$，一元二次方程 $x^2 - ax - 1 = 0$ 有实根；

\vspace{0.3em}
(2) 每个正方形都是平行四边形；

\vspace{0.3em}
(3) $\exists m \in \mathbf{N}$，$\sqrt{m^2 + 1} \in \mathbf{N}$；

\vspace{0.3em}
(4) 存在一个四边形 $ABCD$，其内角和不等于 $360^\circ$。
\end{frame}

\begin{frame}[t]{第7题 解答}
\begin{tcolorbox}[colback=orange!10,colframe=orange,title=\textbf{解}]
(1) 否定：$\exists a \in \mathbf{R}$，方程 $x^2 - ax - 1 = 0$ 无实根。

\vspace{0.3em}
原命题：判别式 $\Delta = a^2 + 4 > 0$ 恒成立，所以原命题\textbf{真}，否定\textbf{假}。

\vspace{0.6em}
(2) 否定：存在一个正方形不是平行四边形。

\vspace{0.3em}
原命题\textbf{真}（正方形是特殊的平行四边形），否定\textbf{假}。

\vspace{0.6em}
(3) 否定：$\forall m \in \mathbf{N}$，$\sqrt{m^2 + 1} \notin \mathbf{N}$。

\vspace{0.3em}
取 $m = 0$，则 $\sqrt{0^2 + 1} = 1 \in \mathbf{N}$，所以原命题\textbf{真}，否定\textbf{假}。

\vspace{0.6em}
(4) 否定：任意一个四边形 $ABCD$，其内角和都等于 $360^\circ$。

\vspace{0.3em}
原命题\textbf{假}（任意四边形内角和都是 $360^\circ$），否定\textbf{真}。
\end{tcolorbox}
\end{frame}

% ========== 第8题 ==========
\begin{frame}[t]{第8题：点集的交集及几何意义}
\textbf{8.} 已知集合 $A = \{(x, y) \mid 2x - y = 0\}$，$B = \{(x, y) \mid 3x + y = 0\}$，$C = \{(x, y) \mid 2x - y = 3\}$，求 $A \cap B$，$A \cap C$，并解释它们的几何意义。
\end{frame}

\begin{frame}[t]{第8题 解答}
\begin{tcolorbox}[colback=orange!10,colframe=orange,title=\textbf{解}]
\textbf{求 $A \cap B$：} 解方程组
\[
\begin{cases}
2x - y = 0, \\
3x + y = 0.
\end{cases}
\]
两式相加得 $5x = 0$，所以 $x = 0$，代入得 $y = 0$。
\[
A \cap B = \{(0,\ 0)\}.
\]
几何意义：直线 $2x - y = 0$ 与直线 $3x + y = 0$ 的交点为原点。

\vspace{0.8em}
\textbf{求 $A \cap C$：} 解方程组
\[
\begin{cases}
2x - y = 0, \\
2x - y = 3.
\end{cases}
\]
两个方程左边相同，右边 $0 \neq 3$，方程组无解。
\[
A \cap C = \varnothing.
\]
几何意义：直线 $2x - y = 0$ 与直线 $2x - y = 3$ 斜率相同（都为 $2$），互相平行，没有交点。
\end{tcolorbox}
\end{frame}

% ========== 第9题 ==========
\begin{frame}[t]{第9题：集合相等问题}
\textbf{9.} 已知集合 $A = \{1,\ 3,\ a^2\}$，$B = \{1,\ a + 2\}$，是否存在实数 $a$，使得 $A \cup B = A$？若存在，试求出实数 $a$ 的值；若不存在，请说明理由。
\end{frame}

\begin{frame}[t]{第9题 解答}
\begin{tcolorbox}[colback=orange!10,colframe=orange,title=\textbf{解}]
若 $A \cup B = A$，则 $B \subseteq A$，即 $B$ 中的元素都属于 $A$。

\vspace{0.5em}
因为 $1 \in A$，所以只需 $a + 2 \in A$，即 $a + 2 = 3$ 或 $a + 2 = a^2$。

\vspace{0.8em}
\textbf{情况1：} $a + 2 = 3$，解得 $a = 1$。

\vspace{0.3em}
此时 $A = \{1,\ 3,\ 1\}$，不满足集合元素的互异性，舍去。

\vspace{0.8em}
\textbf{情况2：} $a + 2 = a^2$，即 $a^2 - a - 2 = 0$，解得 $a = 2$ 或 $a = -1$。

\vspace{0.3em}
当 $a = 2$ 时，$A = \{1,\ 3,\ 4\}$，$B = \{1,\ 4\}$，满足 $B \subseteq A$。

\vspace{0.3em}
当 $a = -1$ 时，$A = \{1,\ 3,\ 1\}$，不满足互异性，舍去。

\vspace{0.8em}
\textbf{综上：} 存在实数 $a = 2$，使得 $A \cup B = A$。
\end{tcolorbox}
\end{frame}

% ========== 第10题 ==========
\begin{frame}[t]{第10题：把定理写成含量词的命题}
\textbf{10.} 把下列定理表示的命题写成含有量词的命题：

\vspace{0.5em}
(1) 勾股定理；

\vspace{0.3em}
(2) 三角形内角和定理。
\end{frame}

\begin{frame}[t]{第10题 解答}
\begin{tcolorbox}[colback=orange!10,colframe=orange,title=\textbf{解}]
(1) \textbf{勾股定理：}

\vspace{0.3em}
$\forall$ 直角三角形，若它的两条直角边长为 $a$、$b$，斜边长为 $c$，则
\[
a^2 + b^2 = c^2.
\]

\vspace{0.8em}
(2) \textbf{三角形内角和定理：}

\vspace{0.3em}
$\forall \triangle ABC$，它的三个内角 $A$、$B$、$C$ 满足
\[
A + B + C = 180^\circ.
\]
\end{tcolorbox}
\end{frame}

% ========== 第11题 ==========
\begin{frame}[t]{第11题：容斥原理应用题}
\textbf{11.} 学校举办运动会时，高一（1）班共有 $28$ 名同学参加比赛，有 $15$ 人参加游泳比赛，有 $8$ 人参加田径比赛，有 $14$ 人参加球类比赛，同时参加游泳比赛和田径比赛的有 $3$ 人，同时参加游泳比赛和球类比赛的有 $3$ 人，没有人同时参加三项比赛。问：同时参加田径和球类比赛的有多少人？只参加游泳一项比赛的有多少人？
\end{frame}

\begin{frame}[t]{第11题 解答}
\begin{tcolorbox}[colback=orange!10,colframe=orange,title=\textbf{解}]
设参加游泳、田径、球类比赛的同学集合分别为 $S$、$T$、$B$。

\vspace{0.5em}
已知：$|S \cup T \cup B| = 28$，$|S| = 15$，$|T| = 8$，$|B| = 14$，$|S \cap T| = 3$，$|S \cap B| = 3$，$|S \cap T \cap B| = 0$。

\vspace{0.5em}
由三集合容斥原理：
\[
|S \cup T \cup B| = |S| + |T| + |B| - |S \cap T| - |S \cap B| - |T \cap B| + |S \cap T \cap B|.
\]

\vspace{0.3em}
代入数据：
\[
28 = 15 + 8 + 14 - 3 - 3 - |T \cap B| + 0.
\]

\vspace{0.3em}
解得：$|T \cap B| = 3$。

\vspace{0.5em}
\textbf{只参加游泳一项比赛的人数：}
\[
|S| - |S \cap T| - |S \cap B| = 15 - 3 - 3 = 9.
\]

\vspace{0.5em}
\textbf{答案：} 同时参加田径和球类比赛的有 $3$ 人，只参加游泳一项比赛的有 $9$ 人。
\end{tcolorbox}
\end{frame}

% ========== 第12题 ==========
\begin{frame}[t]{第12题：写出含量词的命题}
\textbf{12.} 根据下述事实，分别写出含有量词的全称量词命题或存在量词命题：

\vspace{0.5em}
(1) 
\[
\begin{aligned}
1 &= 1^2, \\
1 + 3 &= 2^2, \\
1 + 3 + 5 &= 3^2, \\
1 + 3 + 5 + 7 &= 4^2, \\
1 + 3 + 5 + 7 + 9 &= 5^2, \\
&\cdots
\end{aligned}
\]

\vspace{0.5em}
(2) 如图，在 $\triangle ABC$ 中，$AD$，$BE$ 与 $CF$ 分别为 $BC$，$AC$ 与 $AB$ 边上的高，则 $AD$，$BE$ 与 $CF$ 所在的直线交于一点 $O$。
\end{frame}

\begin{frame}[t]{第12题 解答}
\begin{tcolorbox}[colback=orange!10,colframe=orange,title=\textbf{解}]
(1) 观察规律：前 $n$ 个连续奇数的和等于 $n^2$。

\vspace{0.3em}
全称量词命题：
\[
\forall n \in \mathbf{N}^*,\quad 1 + 3 + 5 + \cdots + (2n-1) = n^2.
\]

\vspace{0.8em}
(2) 全称量词命题：

\vspace{0.3em}
$\forall \triangle ABC$，它的三条高所在的直线交于一点（即三角形的垂心）。
\end{tcolorbox}
\end{frame}

\begin{frame}[c]
  \begin{center}
    {\Huge\calligra Thanks!}
  \end{center}
\end{frame}

\end{document}
